% Options for packages loaded elsewhere
\PassOptionsToPackage{unicode}{hyperref}
\PassOptionsToPackage{hyphens}{url}
\PassOptionsToPackage{dvipsnames,svgnames,x11names}{xcolor}
\documentclass[
  10pt,
  fontset=fandol]{ctexart}
\usepackage{xcolor}
\usepackage[margin=16mm]{geometry}
\usepackage{amsmath,amssymb}
\setcounter{secnumdepth}{-\maxdimen} % remove section numbering
\usepackage{iftex}
\ifPDFTeX
  \usepackage[T1]{fontenc}
  \usepackage[utf8]{inputenc}
  \usepackage{textcomp} % provide euro and other symbols
\else % if luatex or xetex
  \usepackage{unicode-math} % this also loads fontspec
  \defaultfontfeatures{Scale=MatchLowercase}
  \defaultfontfeatures[\rmfamily]{Ligatures=TeX,Scale=1}
\fi
\usepackage{lmodern}
\ifPDFTeX\else
  % xetex/luatex font selection
\fi
% Use upquote if available, for straight quotes in verbatim environments
\IfFileExists{upquote.sty}{\usepackage{upquote}}{}
\IfFileExists{microtype.sty}{% use microtype if available
  \usepackage[]{microtype}
  \UseMicrotypeSet[protrusion]{basicmath} % disable protrusion for tt fonts
}{}
\makeatletter
\@ifundefined{KOMAClassName}{% if non-KOMA class
  \IfFileExists{parskip.sty}{%
    \usepackage{parskip}
  }{% else
    \setlength{\parindent}{0pt}
    \setlength{\parskip}{6pt plus 2pt minus 1pt}}
}{% if KOMA class
  \KOMAoptions{parskip=half}}
\makeatother
\setlength{\emergencystretch}{3em} % prevent overfull lines
\providecommand{\tightlist}{%
  \setlength{\itemsep}{0pt}\setlength{\parskip}{0pt}}
\usepackage{amsmath,amssymb,mathtools,needspace}
\xeCJKDeclareCharClass{CJK}{"2032,"0400->"04FF,"2160->"216B,"2460->"2473,"25A0,"25C7,"25CB,"25CF}
\IfFontExistsTF{Arial Unicode MS}{\xeCJKsetup{AutoFallBack=true}\setCJKfallbackfamilyfont{\CJKrmdefault}{Arial Unicode MS}}{}
\setcounter{secnumdepth}{0}
\setlength{\emergencystretch}{2em}
\allowdisplaybreaks
\usepackage{bookmark}
\IfFileExists{xurl.sty}{\usepackage{xurl}}{} % add URL line breaks if available
\urlstyle{same}
\hypersetup{
  pdftitle={梯形法的余项展开式},
  colorlinks=true,
  linkcolor={Maroon},
  filecolor={Maroon},
  citecolor={Blue},
  urlcolor={Blue},
  pdfcreator={LaTeX via pandoc}}

\title{梯形法的余项展开式}
\author{}
\date{}

\begin{document}
\maketitle

\subsubsection{4.3.4 梯形法的余项展开式}\label{ux68afux5f62ux6cd5ux7684ux4f59ux9879ux5c55ux5f00ux5f0f}

前述外推方法的基础是梯形法的余项展开式（4.3.7），现在运用 Taylor 展开方法推导这个公式。

将 \(f(x)\) 在子区间 \([x_k,x_{k+1}]\) 的中点 \(x_{k+\frac12}=\dfrac12(x_k+x_{k+1})\) 展开，有

\[
\begin{aligned}
f(x)={}&f_{k+\frac12}+(x-x_{k+\frac12})f'_{k+\frac12}
+\frac{(x-x_{k+\frac12})^2}{2!}f''_{k+\frac12}
+\frac{(x-x_{k+\frac12})^3}{3!}f'''_{k+\frac12}\\
&+\frac{(x-x_{k+\frac12})^4}{4!}f^{(4)}_{k+\frac12}
+\frac{(x-x_{k+\frac12})^5}{5!}f^{(5)}_{k+\frac12}
+\frac{(x-x_{k+\frac12})^6}{6!}f^{(6)}_{k+\frac12}+\cdots,
\end{aligned}
\tag{4.3.14}
\]

其中 \(f^{(j)}_{k+\frac12}\) 是 \(f^{(j)}(x_{k+\frac12})\) 的缩写。据此写出 \(f(x_k)\) 和 \(f(x_{k+1})\) 的展开式，易知

\[
\frac h2[f(x_k)+f(x_{k+1})]
=hf_{k+\frac12}+\frac h{2!}\left(\frac h2\right)^2f''_{k+\frac12}
+\frac h{4!}\left(\frac h2\right)^4f^{(4)}_{k+\frac12}
+\frac h{6!}\left(\frac h2\right)^6f^{(6)}_{k+\frac12}+\cdots,
\]

求和得

\[
\begin{aligned}
T(h)&=\frac h2\sum[f(x_k)+f(x_{k+1})]\\
&=h\sum f_{k+\frac12}+\frac{h^3}{2!\times2^2}\sum f''_{k+\frac12}
+\frac{h^5}{4!\times2^4}\sum f^{(4)}_{k+\frac12}
+\frac{h^7}{6!\times2^6}\sum f^{(6)}_{k+\frac12}+\cdots.
\end{aligned}
\tag{4.3.15}
\]

\href{../../source-images/pdf-106.jpeg}{原页图像}

\subsubsection{4.3.4 梯形法的余项展开式（续）}\label{ux68afux5f62ux6cd5ux7684ux4f59ux9879ux5c55ux5f00ux5f0fux7eed}

为简洁起见，这里省略了符号 \(\displaystyle\sum_{k=0}^{n-1}\) 中的上、下限。

另一方面，将式（4.3.14）在子区间 \([x_k,x_{k+1}]\) 上求积分，有

\[
\int_{x_k}^{x_{k+1}}f(x)\,\mathrm{d}x
=hf_{k+\frac12}
+\frac{2}{3!}\left(\frac h2\right)^3f''_{k+\frac12}
+\frac{2}{5!}\left(\frac h2\right)^5f^{(4)}_{k+\frac12}
+\frac{2}{7!}\left(\frac h2\right)^7f^{(6)}_{k+\frac12}+\cdots,
\]

然后再关于 \(k\) 从 \(0\) 到 \(n-1\) 求和，得

\[
\begin{aligned}
I&=\int_a^b f(x)\,\mathrm{d}x\\
&=h\sum f_{k+\frac12}
+\frac{h^3}{3!\times2^2}\sum f''_{k+\frac12}
+\frac{h^5}{5!\times2^4}\sum f^{(4)}_{k+\frac12}
+\frac{h^7}{7!\times2^6}\sum f^{(6)}_{k+\frac12}+\cdots.
\end{aligned}
\tag{4.3.16}
\]

利用式（4.3.16）从式（4.3.15）中消去项 \(h\sum f_{k+\frac12}\)，得

\[
T(h)=I+\frac{h^2}{2!\times6}\sum f''_{k+\frac12}
+\frac{h^5}{4!\times20}\sum f^{(4)}_{k+\frac12}
+\frac{3h^7}{6!\times224}\sum f^{(6)}_{k+\frac12}+\cdots.
\tag{4.3.17}
\]

又对 \(f''(x)\) 应用式（4.3.16），并注意到

\[
\int_a^b f''(x)\,\mathrm{d}x=f'(b)-f'(a),
\]

有

\[
h\sum f''_{k+\frac12}=f'(b)-f'(a)
-\frac{h^3}{3!\times2^2}\sum f^{(4)}_{k+\frac12}
-\frac{h^5}{5!\times2^4}\sum f^{(6)}_{k+\frac12}+\cdots,
\]

代入式（4.3.17），整理得

\[
T(h)=I+\frac{h^2}{2!\times6}[f'(b)-f'(a)]
-\frac{h^5}{4!\times30}\sum f^{(4)}_{k+\frac12}
-\frac{h^7}{6!\times56}\sum f^{(6)}_{k+\frac12}+\cdots.
\]

再对 \(f^{(4)}(x)\) 应用式（4.3.16），有

\[
h\sum f^{(4)}_{k+\frac12}=f'''(b)-f'''(a)
-\frac{h^3}{3!\times2^2}\sum f^{(6)}_{k+\frac12}+\cdots,
\]

从而可进一步得

\[
\begin{aligned}
T(h)={}&I+\frac{h^2}{2!\times6}[f'(b)-f'(a)]
-\frac{h^4}{4!\times30}[f'''(b)-f'''(a)]\\
&+\frac{h^7}{6!\times42}\sum f^{(6)}_{k+\frac12}-\cdots.
\end{aligned}
\tag{4.3.18}
\]

反复施行上述手续，即可得到形如式（4.3.7）的余项公式。

最后再强调一下，应用 Romberg 方法时，一定要注意余项展开式（4.3.7）成立这个前提，不然就得不出正确的结果。

\begin{center}\rule{0.5\linewidth}{0.5pt}\end{center}

\subsection{相关原页校注（agent 补充）}\label{ux76f8ux5173ux539fux9875ux6821ux6ce8agent-ux8865ux5145}

以下校注来自本节覆盖的原页，可能也涉及同页相邻小节；原文照录与校正说明分开保留。

\subsubsection{原书 PDF 页 106 的校注}\label{ux539fux4e66-pdf-ux9875-106-ux7684ux6821ux6ce8}

\href{../pages/pdf-106.md}{完整原页转录} · \href{../../source-images/pdf-106.jpeg}{原页图像}

校注记录：CH04B-E001。

\begin{quote}
校注（agent 补充，CH04B-E001）：原扫描式（4.3.17）首个分子明确印为 \(h^2\)，本转录原样保留；该处疑为原书排印错误。由本页式（4.3.16）及中点展开的梯形项相减，\(f''\) 项系数应为 \(h^3/12=h^3/(2!\times6)\)，才能在代入紧接着的 \(h\sum f''\) 关系后得到原页下一式的 \(h^2[f'(b)-f'(a)]/12\)。这不是 OCR 改写；\href{../assets/ch04b/review-p106-middle.png}{局部原图证据}。
\end{quote}

\subsection{本节覆盖原页的校注记录}\label{ux672cux8282ux8986ux76d6ux539fux9875ux7684ux6821ux6ce8ux8bb0ux5f55}

下列记录按原页关联，含同页相邻小节的校注；计算时同时读取上文校注和完整原页。

\begin{itemize}
\tightlist
\item
  \texttt{CH04B-E001}：原书 PDF 页 106
\end{itemize}

\end{document}
